	\section{Error handling}

	All isotherm models should be fitted to experimental data using
	\emph{non-linear} regression (e.g.\ Levenberg--Marquardt or
	trust-region algorithms) applied directly to the original equation
	$q_e = f(C_e;\,\boldsymbol{\theta})$, where
	$\boldsymbol{\theta}$ is the parameter vector.
	Linearised forms (e.g.\ $1/q_e$ vs.\ $1/C_e$ for Langmuir) should
	be avoided because the transformation distorts the error structure
	and biases the estimated parameters.

	\subsubsection{Residual error metrics}

	Let $n$ be the number of experimental points,
	$q_{e,i}^{\exp}$ the measured value,
	$q_{e,i}^{\mathrm{pred}}$ the model prediction, and $p$ the
	number of fitted parameters. The following metrics quantify the
	goodness of fit.

	\paragraph{Sum of Squared Errors (SSE).}
	The SSE is the most fundamental measure of total discrepancy
	between model and data. It sums the squared residuals, thereby
	giving greater weight to large deviations. Because it is
	unbounded and depends on the units and magnitude of $q_e$, SSE
	values can only be compared across models fitted to the
	\emph{same} dataset. It is the objective function minimised in
	ordinary least-squares regression:
	\begin{equation}
		\mathrm{SSE}
		= \sum_{i=1}^{n}
		\bigl(q_{e,i}^{\exp} - q_{e,i}^{\mathrm{pred}}\bigr)^{2}.
		\label{eq:SSE}
	\end{equation}
	A smaller SSE indicates a better fit. However, SSE always
	decreases (or stays the same) when parameters are added, so it
	cannot by itself distinguish between a genuinely better model and
	overfitting.

	\paragraph{Coefficient of determination ($R^{2}$).}
	$R^{2}$ normalises the SSE by the total sum of squares (SST),
	giving a dimensionless measure of the fraction of variance in the
	data that is explained by the model:
	\begin{equation}
		R^{2} = 1 - \frac{\mathrm{SSE}}{\mathrm{SST}},
		\qquad
		\mathrm{SST} = \sum_{i=1}^{n}
		\bigl(q_{e,i}^{\exp} - \bar{q}_e^{\exp}\bigr)^{2},
		\label{eq:R2}
	\end{equation}
	where $\bar{q}_e^{\exp}$ is the mean of the experimental values.
	$R^{2} = 1$ means a perfect fit; $R^{2} = 0$ means the model
	is no better than a horizontal line at $\bar{q}_e$. Although
	widely reported, $R^{2}$ has important limitations: (i)~it cannot
	penalise over-parameterisation (a model with $p = n$ always
	gives $R^{2} = 1$); (ii)~for non-linear models, $R^{2}$ is not
	guaranteed to lie between 0 and 1; (iii)~a high $R^{2}$ does not
	guarantee that the model captures the correct physics.

	\paragraph{Adjusted $R^{2}$.}
	The adjusted coefficient of determination corrects $R^{2}$ for
	the number of fitted parameters, penalising unnecessary complexity.
	When an additional parameter does not improve the fit sufficiently,
	$R^{2}_{\mathrm{adj}}$ can actually \emph{decrease}, unlike
	$R^{2}$:
	\begin{equation}
		R_{\mathrm{adj}}^{2}
		= 1 - \frac{n - 1}{n - p - 1}\,(1 - R^{2}).
		\label{eq:R2adj}
	\end{equation}
	The factor $(n-1)/(n-p-1)$ inflates the residual variance when
	$p$ is large relative to $n$, making $R^{2}_{\mathrm{adj}}$ a
	more honest indicator than $R^{2}$ in model comparison.

	\paragraph{Sum of Absolute Errors (EABS).}
	EABS replaces the squared residuals of SSE with absolute values,
	making it less sensitive to outliers. A single anomalous data
	point that would dominate the SSE contributes only linearly to
	EABS:
	\begin{equation}
		\mathrm{EABS}
		= \sum_{i=1}^{n}
		\bigl|q_{e,i}^{\exp} - q_{e,i}^{\mathrm{pred}}\bigr|.
		\label{eq:EABS}
	\end{equation}
	Like SSE, EABS is scale-dependent and should only be compared
	across models fitted to the same dataset.

	\paragraph{Mean Absolute Residual (RESID).}
	RESID is simply the EABS divided by $n$, giving the average
	absolute prediction error in the same units as $q_e$. It
	provides an intuitive measure of ``how wrong, on average, is
	the model prediction'':
	\begin{equation}
		\mathrm{RESID}
		= \frac{1}{n}\sum_{i=1}^{n}
		\bigl|q_{e,i}^{\exp} - q_{e,i}^{\mathrm{pred}}\bigr|.
		\label{eq:RESID}
	\end{equation}
	For example, $\mathrm{RESID} = 0.5$~mg/g means the model
	predictions deviate from the data by 0.5~mg/g on average.

	\paragraph{Average Relative Error Deviation (ARED).}
	ARED expresses the error as a percentage of the experimental
	value, making it dimensionless and comparable across datasets
	with different $q_e$ magnitudes. It is particularly useful for
	adsorption data where $q_e$ spans several orders of magnitude,
	because it gives equal weight to low-loading and high-loading
	points:
	\begin{equation}
		\mathrm{ARED}
		= \frac{100}{n}\sum_{i=1}^{n}
		\left|
		\frac{q_{e,i}^{\exp} - q_{e,i}^{\mathrm{pred}}}
		{q_{e,i}^{\exp}}
		\right|.
		\label{eq:ARED}
	\end{equation}
	Caution: ARED can become very large (or undefined) if
	$q_{e,i}^{\exp}$ is close to zero; data points with very small
	$q_e$ should be checked for excessive influence on ARED.

	\paragraph{Marquardt's Percent Standard Deviation (MPSED).}
	MPSED uses the \emph{median} (instead of the mean) of the
	squared relative errors, making it a robust statistic that is
	largely immune to outliers. A single bad data point that
	inflates ARED will have minimal effect on MPSED:
	\begin{equation}
		\mathrm{MPSED}
		= 100 \times \mathrm{median}\!\left[
		\left(
		\frac{q_{e,i}^{\exp} - q_{e,i}^{\mathrm{pred}}}
		{q_{e,i}^{\exp}}
		\right)^{\!2}
		\right].
		\label{eq:MPSED}
	\end{equation}
	MPSED is especially valuable when the dataset contains a few
	points with large experimental uncertainty (e.g.\ near the
	detection limit).

	\paragraph{Hybrid Fractional Error Function (HYBRID).}
	The HYBRID function combines features of SSE and ARED: it sums
	squared residuals weighted by $1/q_{e,i}^{\exp}$, giving a
	balance between absolute and relative errors. The normalisation
	by the degrees of freedom $(n-p)$ instead of $n$ penalises
	models with more parameters:
	\begin{equation}
		\mathrm{HYBRID}
		= \frac{100}{n - p}\sum_{i=1}^{n}
		\frac{\bigl(q_{e,i}^{\exp} - q_{e,i}^{\mathrm{pred}}\bigr)^{2}}
		{q_{e,i}^{\exp}}.
		\label{eq:HYBRID}
	\end{equation}
	Because the denominator is $q_{e,i}^{\exp}$ (not its square),
	the HYBRID function gives intermediate sensitivity between SSE
	(which favours fitting high-$q_e$ points) and ARED (which
	favours fitting low-$q_e$ points). The built-in degrees-of-freedom
	correction makes HYBRID one of the most balanced error functions
	for isotherm model comparison.

	\subsubsection{Information criteria for model selection}

	When comparing models with different numbers of parameters, pure
	error metrics (SSE, $R^{2}$) are insufficient because a model with
	more parameters will always fit better. Information criteria
	penalise model complexity and are the correct tool for model
	selection.

	\paragraph{Akaike Information Criterion (AIC).}
	\begin{equation}
		\mathrm{AIC}
		= n\,\ln\!\left(\frac{\mathrm{SSE}}{n}\right) + 2p.
		\label{eq:AIC}
	\end{equation}
	For small sample sizes ($n/p < 40$), the corrected version should
	be used:
	\begin{equation}
		\mathrm{AIC}_c
		= \mathrm{AIC} + \frac{2p(p+1)}{n - p - 1}.
		\label{eq:AICc}
	\end{equation}

	\paragraph{Bayesian Information Criterion (BIC).}
	Applies a stronger penalty for additional parameters:
	\begin{equation}
		\mathrm{BIC}
		= n\,\ln\!\left(\frac{\mathrm{SSE}}{n}\right)
		+ p\,\ln n.
		\label{eq:BIC}
	\end{equation}

	In both cases, the model with the \emph{lowest} AIC (or BIC) is
	preferred. The difference $\Delta\mathrm{AIC}_i = \mathrm{AIC}_i
	- \mathrm{AIC}_{\min}$ can be used to rank models:
	\begin{itemize}
		\item $\Delta\mathrm{AIC} < 2$: substantial support for
		the model.
		\item $4 < \Delta\mathrm{AIC} < 7$: considerably less support.
		\item $\Delta\mathrm{AIC} > 10$: essentially no support.
	\end{itemize}

	\subsubsection{F-test for nested models}

	When a simpler model (model~1, $p_1$ parameters) is a special case
	of a more complex model (model~2, $p_2 > p_1$ parameters), the
	F-test determines whether the improvement in fit justifies the
	additional parameters:
	\begin{equation}
		F = \frac{(\mathrm{SSE}_1 - \mathrm{SSE}_2)\,/\,
			(p_2 - p_1)}
		{\mathrm{SSE}_2\,/\,(n - p_2)}.
		\label{eq:Ftest}
	\end{equation}
	Under the null hypothesis (simpler model is adequate), $F$ follows
	an $F$-distribution with $(p_2 - p_1,\; n - p_2)$ degrees of
	freedom. If $F > F_{\alpha}$, the more complex model is
	statistically justified at significance level $\alpha$.

	Example: Langmuir ($p_1 = 2$) vs.\ Sips ($p_2 = 3$), or
	Langmuir ($p_1 = 2$) vs.\ Double Langmuir ($p_2 = 4$).

	\subsubsection{Practical recommendations}

	\begin{enumerate}
		\item Always report \emph{at least} $R^{2}$, SSE, and one
		relative metric (ARED or HYBRID).
		\item When comparing models with different $p$, use AIC$_c$
		or BIC rather than $R^{2}$.
		\item Use the F-test for nested model pairs
		(Langmuir vs.\ Double Langmuir, Langmuir vs.\ Sips).
		\item Visualise residuals: non-random patterns indicate
		systematic model inadequacy regardless of the $R^{2}$ value.
		\item Report confidence intervals on fitted parameters
		(from the Jacobian or bootstrap), not just point estimates.
		\item For heatmap visualisation of error metrics across
		materials and models, apply a $\log_{10}$ transformation to
		compress the dynamic range.
	\end{enumerate}

